\documentclass[11pt]{article}
\usepackage[margin=1in]{geometry}
\usepackage{amsmath, amssymb}
\usepackage{enumitem}
\usepackage{booktabs}
\usepackage{fancyhdr}
\usepackage{minted}
\usepackage{xcolor}

\pagestyle{fancy}
\fancyhf{}
\lhead{EC 282: Introduction to Econometrics}
\rhead{Spring 2026}
\cfoot{\thepage}

\setminted{
  bgcolor=gray!5,
  fontsize=\small,
  linenos,
  frame=single,
  framesep=2mm,
  breaklines
}

\begin{document}

\begin{center}
{\LARGE \textbf{Handout 6: Inference and Binary Regressors}} \\[8pt]
{\large EC 282: Introduction to Econometrics} \\[4pt]
Spring 2026
\end{center}

\vspace{0.3cm}

\noindent \textbf{Instructions:} Run the provided R code and answer the questions. Show your work for calculations. This handout builds directly on Handout~5---make sure you have that code loaded first.

\section{Setup: Continuing with the Gapminder Data}

We continue with the same Gapminder 2007 cross-section from Handout~5. Run the code below to reload the data and re-estimate the regression.

\begin{minted}{r}
library(ggplot2)

options(scipen = 999)

# Reload the Gapminder 2007 data
gapminder_full <- read.csv(
  "https://raw.githubusercontent.com/resbaz/r-novice-gapminder-files/
   master/data/gapminder-FiveYearData.csv"
)
gap07 <- gapminder_full[gapminder_full$year == 2007, ]

# Re-estimate the simple regression from Handout 5
reg <- lm(lifeExp ~ gdpPercap, data = gap07)
\end{minted}

\section{Reading the Regression Output}

\textbf{Question 2.1:} Run \texttt{summary(reg)} and examine the output carefully.

\begin{minted}{r}
summary(reg)
\end{minted}

\vspace{0.5cm}

Identify and write down the following from the output:

\begin{enumerate}[label=(\alph*)]
\item The estimated slope $\hat{\beta}_1$ and intercept $\hat{\beta}_0$
\item The standard error of $\hat{\beta}_1$, denoted $SE(\hat{\beta}_1)$
\item The t-statistic for $\hat{\beta}_1$
\item The p-value for $\hat{\beta}_1$
\item The $R^2$
\end{enumerate}

\vspace{2cm}

\textbf{Question 2.2:} The standard error measures the \textit{precision} of our estimate. In your own words, what does a small standard error tell you? What does a large one tell you?

\vspace{3cm}

\textbf{Question 2.3:} What affects the precision of $\hat{\beta}_1$? Three factors determine $SE(\hat{\beta}_1)$: the sample size $n$, the spread of $X$ (i.e., $\text{Var}(X_i)$), and the amount of noise in the regression (i.e., $\text{Var}(u_i)$). For each factor, explain whether \textit{increasing} it makes $SE(\hat{\beta}_1)$ larger or smaller, and why.

\vspace{4cm}

\section{Hypothesis Testing}

We want to test whether GDP per capita has a \textit{statistically significant} relationship with life expectancy. This means testing whether the true slope $\beta_1$ could be zero.

\textbf{Question 3.1:} Write down the null and alternative hypotheses for a two-sided test of whether GDP per capita is related to life expectancy.

\vspace{2cm}

\textbf{Question 3.2:} The t-statistic is computed as:
\[
t\text{-stat} = \frac{\hat{\beta}_1 - \beta_{1,0}}{SE(\hat{\beta}_1)}
\]
where $\beta_{1,0}$ is the hypothesized value under $H_0$. Using the values from the regression output, compute the t-statistic by hand. Verify it matches the output.

\begin{minted}{r}
# Extract coefficients and standard errors
beta1_hat <- coef(reg)["gdpPercap"]
se_beta1  <- summary(reg)$coefficients["gdpPercap", "Std. Error"]

cat("beta1_hat:", beta1_hat, "\n")
cat("SE(beta1_hat):", se_beta1, "\n")

# Compute t-statistic by hand
t_stat <- beta1_hat / se_beta1
cat("t-statistic (by hand):", t_stat, "\n")
\end{minted}

\vspace{2cm}

\textbf{Question 3.3:} To make a decision, we compare the t-statistic to critical values. For a two-sided test at the 5\% significance level, we reject $H_0$ if $|t| > 1.96$.

\begin{enumerate}[label=(\alph*)]
\item Is the slope on GDP per capita statistically significant at the 5\% level?
\item Is it significant at the 1\% level? (Critical value: 2.576)
\end{enumerate}

\vspace{2cm}

\textbf{Question 3.4:} The p-value gives the probability of observing a t-statistic at least as extreme as ours, \textit{assuming $H_0$ is true}. Compute the p-value in R and verify it matches the regression output.

\begin{minted}{r}
# Compute p-value for two-sided test
p_value <- 2 * pnorm(-abs(t_stat))
cat("p-value (by hand):", p_value, "\n")
cat("p-value (from summary):",
    summary(reg)$coefficients["gdpPercap", "Pr(>|t|)"], "\n")
\end{minted}

\vspace{1cm}

Interpret the p-value in a sentence. What does it tell us about the relationship between GDP per capita and life expectancy?

\vspace{2cm}

\textbf{Question 3.5:} Suppose instead we wanted to test whether GDP per capita has a \textit{positive} effect on life expectancy (a one-sided test):
\[
H_0: \beta_1 = 0 \qquad H_A: \beta_1 > 0
\]

\begin{minted}{r}
# One-sided p-value (right tail)
p_value_one <- pnorm(-abs(t_stat))
cat("One-sided p-value:", p_value_one, "\n")
\end{minted}

How does the one-sided p-value compare to the two-sided p-value? Why?

\vspace{2cm}

\section{Confidence Intervals}

A confidence interval provides a range of plausible values for the true population parameter $\beta_1$.

\textbf{Question 4.1:} A 95\% confidence interval for $\beta_1$ is:
\[
\hat{\beta}_1 \pm 1.96 \times SE(\hat{\beta}_1)
\]

Compute this by hand using the values from the regression output.

\begin{minted}{r}
# 95% confidence interval by hand
ci_lower <- beta1_hat - 1.96 * se_beta1
ci_upper <- beta1_hat + 1.96 * se_beta1

cat("95% CI: [", round(ci_lower, 7), ",", round(ci_upper, 7), "]\n")
\end{minted}

\vspace{1cm}

Verify with R's built-in function:

\begin{minted}{r}
confint(reg, "gdpPercap", level = 0.95)
\end{minted}

\vspace{1cm}

\textbf{Question 4.2:} Interpret the 95\% confidence interval in a sentence. Does the interval include zero? What does this tell you?

\vspace{3cm}

\textbf{Question 4.3:} Now construct a 99\% confidence interval. Is it wider or narrower than the 95\% CI? Why?

\begin{minted}{r}
# 99% confidence interval
ci99_lower <- beta1_hat - 2.576 * se_beta1
ci99_upper <- beta1_hat + 2.576 * se_beta1

cat("99% CI: [", round(ci99_lower, 7), ",", round(ci99_upper, 7), "]\n")
confint(reg, "gdpPercap", level = 0.99)
\end{minted}

\vspace{2cm}

\textbf{Question 4.4:} The connection between confidence intervals and hypothesis testing: if the 95\% CI does not include a hypothesized value $\beta_{1,0}$, then we would reject $H_0: \beta_1 = \beta_{1,0}$ at the 5\% level. Explain in your own words why this equivalence holds.

\vspace{3cm}

\textbf{Question 4.5:} Suppose we want a confidence interval for the \textit{predicted change} in life expectancy when GDP per capita increases by \$10,000. Compute a 95\% CI for $\Delta Y = \beta_1 \times 10{,}000$.

\begin{minted}{r}
# CI for predicted change when X increases by $10,000
delta_X <- 10000
pred_change <- beta1_hat * delta_X
ci_change_lower <- ci_lower * delta_X
ci_change_upper <- ci_upper * delta_X

cat("Predicted change:", round(pred_change, 2), "years\n")
cat("95% CI for change: [", round(ci_change_lower, 2), ",",
    round(ci_change_upper, 2), "] years\n")
\end{minted}

\vspace{1cm}

Interpret this result in a sentence.

\vspace{2cm}

\section{Regression with Binary Variables}

So far, our regressor (GDP per capita) has been continuous. Now we explore what happens when the independent variable is \textbf{binary} (taking only values 0 or 1).

\textbf{Question 5.1:} Create a binary variable \texttt{africa} that equals 1 if a country is in Africa and 0 otherwise.

\begin{minted}{r}
# Create binary variable
gap07$africa <- ifelse(gap07$continent == "Africa", 1, 0)

cat("Number of African countries:", sum(gap07$africa), "\n")
cat("Number of non-African countries:", sum(1 - gap07$africa), "\n")
\end{minted}

\vspace{1cm}

\textbf{Question 5.2:} Compute the average life expectancy separately for African and non-African countries.

\begin{minted}{r}
# Group means
mean_africa     <- mean(gap07$lifeExp[gap07$africa == 1])
mean_non_africa <- mean(gap07$lifeExp[gap07$africa == 0])

cat("Mean life exp (Africa):", round(mean_africa, 2), "years\n")
cat("Mean life exp (non-Africa):", round(mean_non_africa, 2), "years\n")
cat("Difference:", round(mean_africa - mean_non_africa, 2), "years\n")
\end{minted}

\vspace{2cm}

\textbf{Question 5.3:} Now run a regression of life expectancy on the \texttt{africa} dummy:
\[
\widehat{\text{LifeExp}}_i = \hat{\beta}_0 + \hat{\beta}_1 \times \text{Africa}_i
\]

\begin{minted}{r}
reg_binary <- lm(lifeExp ~ africa, data = gap07)
summary(reg_binary)
\end{minted}

\begin{enumerate}[label=(\alph*)]
\item What does $\hat{\beta}_0$ represent? Compare it to one of the group means you computed above.
\item What does $\hat{\beta}_1$ represent? Compare it to the difference in group means.
\item Is the difference statistically significant at the 5\% level? How do you know?
\end{enumerate}

\vspace{4cm}

\textbf{Question 5.4:} This is a key insight: \textbf{regression with a binary regressor is the same as a two-sample comparison of means}. Verify this by running a t-test and comparing the results.

\begin{minted}{r}
# Two-sample t-test
t.test(lifeExp ~ africa, data = gap07, var.equal = FALSE)
\end{minted}

Compare the difference in means and p-value from the t-test to the regression output. What do you notice?

\vspace{3cm}

\textbf{Question 5.5:} Construct a 95\% confidence interval for $\hat{\beta}_1$ from the binary regression. What does this interval tell us about the life expectancy gap between African and non-African countries?

\begin{minted}{r}
confint(reg_binary, "africa", level = 0.95)
\end{minted}

\vspace{3cm}

\textbf{Question 5.6:} Create a visualization comparing the two groups.

\begin{minted}{r}
ggplot(gap07, aes(x = factor(africa, labels = c("Non-Africa", "Africa")),
                  y = lifeExp)) +
  geom_boxplot(fill = c("steelblue", "coral"), alpha = 0.7) +
  geom_jitter(width = 0.2, alpha = 0.4, size = 1.5) +
  labs(title = "Life Expectancy: Africa vs. Rest of World (2007)",
       x = "", y = "Life Expectancy (years)") +
  theme_minimal()
\end{minted}

What patterns do you notice in the data? Is there overlap between the two groups?

\vspace{3cm}

\section{Bringing It Together}

\textbf{Question 6.1:} We've seen that the relationship between GDP per capita and life expectancy is not perfectly linear (Handout~5 showed curvature in the residual plot). One concern is \textbf{omitted variable bias}. What variables might be in the error term $u_i$ that are both:
\begin{enumerate}
\item Correlated with GDP per capita, AND
\item Affect life expectancy?
\end{enumerate}

Name at least two such variables and explain why they create bias.

\vspace{4cm}

\textbf{Question 6.2:} Compare the two regressions we've run in this handout. Fill in the table:

\begin{center}
\begin{tabular}{l|c|c}
\toprule
 & GDP per capita regression & Africa dummy regression \\
\midrule
$\hat{\beta}_0$ & & \\[6pt]
$\hat{\beta}_1$ & & \\[6pt]
$SE(\hat{\beta}_1)$ & & \\[6pt]
t-statistic & & \\[6pt]
$R^2$ & & \\
\bottomrule
\end{tabular}
\end{center}

\begin{minted}{r}
# Side-by-side comparison
cat("=== GDP per Capita Regression ===\n")
cat("beta0:", round(coef(reg)[1], 4), "\n")
cat("beta1:", round(coef(reg)[2], 7), "\n")
cat("SE(beta1):", round(summary(reg)$coefficients[2, 2], 7), "\n")
cat("t-stat:", round(summary(reg)$coefficients[2, 3], 3), "\n")
cat("R-squared:", round(summary(reg)$r.squared, 4), "\n\n")

cat("=== Africa Dummy Regression ===\n")
cat("beta0:", round(coef(reg_binary)[1], 4), "\n")
cat("beta1:", round(coef(reg_binary)[2], 4), "\n")
cat("SE(beta1):", round(summary(reg_binary)$coefficients[2, 2], 4), "\n")
cat("t-stat:", round(summary(reg_binary)$coefficients[2, 3], 3), "\n")
cat("R-squared:", round(summary(reg_binary)$r.squared, 4), "\n")
\end{minted}

\vspace{1cm}

Which regression has a higher $R^2$? Does this mean it is a ``better'' model? Discuss.

\vspace{3cm}

\textbf{Question 6.3:} Reflect on what we've learned across Handouts 5 and 6. We can now:
\begin{itemize}
\item Estimate a regression line (Handout 5)
\item Assess how well it fits (Handout 5)
\item Test whether the relationship is statistically significant (Handout 6)
\item Construct confidence intervals for the true effect (Handout 6)
\item Use binary variables to compare group means within a regression framework (Handout 6)
\end{itemize}

What is one important thing we still \textit{cannot} do with simple regression? (Hint: Think about why the coefficient on GDP per capita might not represent a causal effect.)

\vspace{3cm}

\end{document}
